\documentclass[]{article}
\usepackage{lmodern}
\usepackage{amssymb,amsmath}
\usepackage{ifxetex,ifluatex}
\usepackage{fixltx2e} % provides \textsubscript
\ifnum 0\ifxetex 1\fi\ifluatex 1\fi=0 % if pdftex
  \usepackage[T1]{fontenc}
  \usepackage[utf8]{inputenc}
\else % if luatex or xelatex
  \ifxetex
    \usepackage{mathspec}
  \else
    \usepackage{fontspec}
  \fi
  \defaultfontfeatures{Ligatures=TeX,Scale=MatchLowercase}
\fi
% use upquote if available, for straight quotes in verbatim environments
\IfFileExists{upquote.sty}{\usepackage{upquote}}{}
% use microtype if available
\IfFileExists{microtype.sty}{%
\usepackage{microtype}
\UseMicrotypeSet[protrusion]{basicmath} % disable protrusion for tt fonts
}{}
\usepackage[margin=1in]{geometry}
\usepackage{hyperref}
\hypersetup{unicode=true,
            pdftitle={LA\_2},
            pdfauthor={Adam Hayes, Erin Omyer, Jasmine Sanchez, and Richard Park},
            pdfborder={0 0 0},
            breaklinks=true}
\urlstyle{same}  % don't use monospace font for urls
\usepackage{graphicx,grffile}
\makeatletter
\def\maxwidth{\ifdim\Gin@nat@width>\linewidth\linewidth\else\Gin@nat@width\fi}
\def\maxheight{\ifdim\Gin@nat@height>\textheight\textheight\else\Gin@nat@height\fi}
\makeatother
% Scale images if necessary, so that they will not overflow the page
% margins by default, and it is still possible to overwrite the defaults
% using explicit options in \includegraphics[width, height, ...]{}
\setkeys{Gin}{width=\maxwidth,height=\maxheight,keepaspectratio}
\IfFileExists{parskip.sty}{%
\usepackage{parskip}
}{% else
\setlength{\parindent}{0pt}
\setlength{\parskip}{6pt plus 2pt minus 1pt}
}
\setlength{\emergencystretch}{3em}  % prevent overfull lines
\providecommand{\tightlist}{%
  \setlength{\itemsep}{0pt}\setlength{\parskip}{0pt}}
\setcounter{secnumdepth}{0}
% Redefines (sub)paragraphs to behave more like sections
\ifx\paragraph\undefined\else
\let\oldparagraph\paragraph
\renewcommand{\paragraph}[1]{\oldparagraph{#1}\mbox{}}
\fi
\ifx\subparagraph\undefined\else
\let\oldsubparagraph\subparagraph
\renewcommand{\subparagraph}[1]{\oldsubparagraph{#1}\mbox{}}
\fi

%%% Use protect on footnotes to avoid problems with footnotes in titles
\let\rmarkdownfootnote\footnote%
\def\footnote{\protect\rmarkdownfootnote}

%%% Change title format to be more compact
\usepackage{titling}

% Create subtitle command for use in maketitle
\newcommand{\subtitle}[1]{
  \posttitle{
    \begin{center}\large#1\end{center}
    }
}

\setlength{\droptitle}{-2em}

  \title{LA\_2}
    \pretitle{\vspace{\droptitle}\centering\huge}
  \posttitle{\par}
    \author{Adam Hayes, Erin Omyer, Jasmine Sanchez, and Richard Park}
    \preauthor{\centering\large\emph}
  \postauthor{\par}
      \predate{\centering\large\emph}
  \postdate{\par}
    \date{1/24/2019}


\begin{document}
\maketitle

\subsection{TEAM NAME}\label{team-name}

\subsection{About Jasmine}\label{about-jasmine}

\begin{figure}
\centering
\includegraphics{file:///C:/Users/sanch/OneDrive/Pictures/Photo\%20Alex.JPG}
\caption{}
\end{figure}

\begin{figure}
\centering
\includegraphics{file:///C:/Users/sanch/OneDrive/Pictures/Photo\%20Niko.JPG}
\caption{}
\end{figure}

\begin{itemize}
\item
  I would like to know what the difference in Educated voters, versus
  non-Educated voters within the US is, and how that fluctuates
  depending on where the voters live, as well as how much access they
  have to voting?
\item
  Six months after graduating, I hope to have gained a full time
  position as a Data Statistician for a compnay. (Preferrably Arrow
  Electronics). Five Years after that, I hope to have gained enough
  insight and support to have transitioned into a leadership role within
  my team, then later within my division of whatever company I will be a
  part of.
\item
  My greatest career accomplishment will be creating more evenly
  accessible platforms of data collection for varrying envrionments and
  communities where interpretation is clear.
\item
  Given these hopes and goals, I want to learn and understand the
  several ways we, as individuals process data and then
  interpret/communicate that information to other entities.
\item
  Fun Fact: I have 2 Siblings, an older sister who is 21 and a younger
  brother who is 11.
\end{itemize}

\subsection{About Erin}\label{about-erin}

\begin{figure}
\centering
\includegraphics{smiley.jpg}
\caption{Picture taken in Seattle, WA}
\end{figure}

\begin{itemize}
\tightlist
\item
  Non-Statistical Question: How many hours does it take for an
  actor/actress to memorize their lines?
\item
  Roughly six months after graduation, I would love to have a job with
  an organization that I am happy to work for, which allow myself to
  start a pathway into my future. Five years beyond graduation, I would
  hope I would be experiencing a life that is continuously enjoy.
  Whether that involve new friends, places, work experience, etc.
\item
  I hope my greatest career accomplishment would be a career where I am
  continuously looking forward to what I am applying in the
  organization. The greatest career won't be making the most amount of
  money, but ensuring that I am content with my career and life that
  follows.
\item
  In Introduction to Data Science, I hope to learn the basic functions
  of coding in R. This will allow myself to expand further into other
  programming languages and accomplish cooperative communication.
\item
  Lastly, an interest fact about myself is that my right ear is smaller
  than my left ear, so I would be called ``nemo'' when I was little.
\end{itemize}


\end{document}
